\section{从廉价 H\&E 翻译到丰富分子表征}

上一章使用同一类基因数据内部的 mask 和空间上下文。本章回到 translation：能否用临床中廉价、普遍的 H\&E 图像，预测昂贵且稀缺的空间转录组或蛋白测量？若成功，模型可把历史病理数据转为更丰富的患者表示；若失败，也可能产生极具迷惑性的伪分子图。

\subsection{多模态补全的训练任务}

Imputation（补全）指根据已观测变量估计缺失变量。它不是简单复制平均值，而是学习条件分布，并且应明确预测不确定性。这里的缺失模态是昂贵 molecular assay。下面从成本轴、条件架构和整模态 mask 三个角度检查：模型是否真的模拟了临床部署时只有 H\&E 的信息条件。

\lecturefigure{slide-213-multimodal-imputation.jpg}{用廉价、普遍的 H\&E 推断昂贵、稀有的分子模态}{00:43:48}

图中成本轴揭示任务价值：H\&E 易获得，空间分子测量昂贵。若模型只在配对数据上训练，却部署到不同医院、不同染色流程或不同肿瘤类型，translation 可能崩溃。因此价值来自跨域泛化与校准，而不只是配对样本重建。

\lecturefigure{slide-215-he-conditioned-model.jpg}{H\&E 编码器条件化 masked gene-count 模型}{00:44:27}

架构把 H\&E patch 编码成条件，再与基因 token 主干结合。H\&E-specific encoder 负责视觉特征，主干负责 count prediction。这个分工允许同一个基因模型在有无图像时运行，但也可能让视觉编码器学习染色批次捷径。

\lecturefigure{slide-218-cross-modal-translation.jpg}{把基因输入全部 mask，迫使模型仅凭 H\&E 预测分子读数}{00:45:00}

当所有 gene tokens 都被 mask，输出完全依赖 H\&E 条件。此时任务从 disambiguation 变成更纯粹的 translation。若仍保留少量真实基因读数，模型则在 translation 与 disambiguation 之间移动：图像给出形态先验，已测 marker 修正不确定性。

\begin{importantbox}{部署时可获得什么，训练目标就应模拟什么}
如果临床部署只有 H\&E，评估时必须把分子输入全部隐藏；若部署允许少量 marker，就应测不同 marker 预算下的性能曲线。训练时偷看完整分子模态会制造不可部署的上限。
\end{importantbox}

\subsection{预测图与误差结构}

最终结果不能只看一张颜色相似的热图。应比较真实测量、插值 baseline、模型预测，并按基因、患者、区域和细胞类型分解误差。进一步还要问：形态上不可识别的分子状态是否对应更宽的不确定区间，模型是否在跨医院与跨染色批次时仍保持校准。

\lecturefigure{slide-224-gene-expression-prediction.jpg}{从 H\&E 预测细胞类型与 KRT19 等基因的空间表达}{00:46:28}

左侧是 H\&E，中间是细胞类型或 overlay，右侧是预测表达。首先比较大尺度结构是否一致，再检查边界、稀有区域和局部峰值。形态与表达相关时模型可能表现很好，但同形异质的分子状态仍可能无法从 H\&E 确定。

\lecturefigure{slide-225-imputation-accuracy.jpg}{不同基因的补全难度显著不同}{00:47:16}

三行对比真实检测、插值和模型预测。某些基因具有清晰形态对应，预测较好；另一些基因几乎不可由 H\&E 唯一确定。平均指标会掩盖这种差异，因此需要 gene-wise calibration：模型应知道自己对哪些目标缺乏可识别信息。

若真实值为 $y$、预测分布均值为 $\hat y$、标准差为 $\hat\sigma$，校准不只要求误差小，还要求
\[
\Pr\left(|y-\hat y|\leq 1.96\hat\sigma\right)\approx 0.95.
\]
区间覆盖率应在 held-out patient、不同中心和不同基因上分别报告。

\begin{warningbox}{形态可预测性有上限}
若两个组织区域在 H\&E 上近似相同，却有不同 RNA 状态，任何只看 H\&E 的模型都无法可靠区分。此时正确输出应是高不确定性，而不是用训练集平均模式生成一张平滑热图。
\end{warningbox}

\lecturefigure{slide-234-representation-pipeline.jpg}{从 H\&E 经模型补全丰富表示，再连接治疗假设}{00:50:01}

这张综合图把 translation 的工程价值说清：先从廉价 H\&E 构造高维患者表征，再把表征用于检索、分层或药物假设。中间任何一步都可能放大偏差。模型补全不是实际测量的无损替代；它更像一个可扩展、带不确定性的虚拟 assay。

\teachervoice{00:43:46--00:48:40，老师强调不同基因的预测质量并不一致，并展示真实检测、插值和模型输出。课堂结论不是“从 H\&E 恢复一切”，而是找出哪些分子信息能被形态可靠约束。}

\begin{knowledgebox}{第一次出现：Calibration}
Calibration（校准）要求置信度与真实正确率匹配。模型若对不可辨认基因仍给出窄区间，就算平均误差尚可，也会误导实验排序。AI for Science 中“不知道自己不知道”通常比偶尔的大误差更危险。
\end{knowledgebox}

\subsection{本章小结}

跨模态 imputation 把 H\&E 的规模优势转化为分子假设，但其上限由可识别信息决定。可靠系统必须在患者级切分上比较真实测量、简单 baseline 和模型预测，按基因报告误差与校准，并把补全结果视为虚拟 assay 而非真实实验。

